\section{نظرة عامة على المنظومة \textenglish{(Application Overview)}}
تُعد منظومة "المقص الطائر" \textenglish{(Flying Shear)} من التطبيقات الكلاسيكية والمتقدمة في آنٍ واحد داخل قطاع الأتمتة الصناعية، وتُستخدم بشكل واسع النطاق في خطوط الإنتاج المستمرة التي تعتمد على السحب أو التدفق المستمر، مثل خطوط إنتاج الصابون، الأنابيب البلاستيكية، والألواح المعدنية.

\subsection{المشكلة الصناعية والحل الهندسي}
في خطوط الإنتاج التقليدية، تتطلب عملية قص المنتج إيقاف السير الناقل تماماً لإجراء القطع، مما يؤدي إلى هدر كبير في الوقت، وانخفاض حاد في الإنتاجية، بالإضافة إلى احتمالية تراكم المواد أو تشوهها أثناء التوقف. من هنا برزت الحاجة لمنظومة "المقص الطائر"، والتي تكمن فكرتها الهندسية الأساسية في تحقيق \textbf{التزامن الحركي الكامل} \textenglish{(Kinematic Synchronization)}. 

حيث تتحرك عربة القص بشكل موازٍ للمنتج وبنفس سرعته الخطية تماماً، مما يجعل السرعة النسبية بينهما تساوي صفراً أثناء فترة القص. هذا التزامن يسمح بقطع المنتج أثناء حركته بدقة متناهية دون الحاجة إلى إيقاف خط الإنتاج، مما يرفع من كفاءة الآلة \textenglish{(OEE)} بشكل جذري.

\subsection{المكونات الأساسية للمنظومة المعمارية}
تتألف معمارية التحكم في هذه المنظومة من عدة محاور ميكانيكية وكهربائية تعمل بتناغم تام:
\begin{itemize}
    \item \textbf{المحور القائد \textenglish{(Master Axis)}:} وهو السير الناقل الذي يحمل شريط الصابون المستمر. يتم مراقبة سرعته وموقعه بدقة متناهية باستخدام مشفر دوراني نبضي \textenglish{(Rotary Encoder)} متصل ميكانيكياً بمحور السير.
    \item \textbf{المحور التابع \textenglish{(Slave Axis)}:} وهو عربة المقص التي تتحرك للأمام والخلف. يتم تحريكها بواسطة محرك سيرفو دقيق \textenglish{(Servo Motor)} قادر على تحقيق تسارع عالٍ واستجابة لحظية للأوامر.
    \item \textbf{العقل المدبر \textenglish{(Controller)}:} يُستخدم في هذا التطبيق متحكم منطقي مبرمج \textenglish{(PLC)} من عائلة \textenglish{Delta DVP}. يتميز هذا المتحكم بقدرته العالية على معالجة العدادات السريعة ذات الـ 32-بت \textenglish{(High-Speed Counters)} لقراءة نبضات المشفر، وقدرته على توليد نبضات عالية التردد \textenglish{(PTO)} للتحكم في درايف السيرفو.
    \item \textbf{آلية القص \textenglish{(Cutting Mechanism)}:} وتتمثل في أسطوانة هوائية \textenglish{(Pneumatic Cylinder)} مسؤولة عن إنزال ورفع شفرة القص في اللحظة التي يتحقق فيها التزامن الحركي التام.
\end{itemize}

\subsection{دورة عمل المنظومة \textenglish{(Operation Cycle)}}
تمر عربة المقص خلال دورة الإنتاج الواحدة بأربع مراحل فيزيائية أساسية، وهي التي سيتم بناء خوارزمية التحكم برمجياً على أساسها:
\begin{enumerate}
    \item \textbf{الانتظار \textenglish{(Standby)}:} وقوف العربة بوضع السكون عند نقطة الصفر بانتظار وصول طول المنتج للحد المطلوب.
    \item \textbf{التسارع والمزامنة \textenglish{(Acceleration \& Sync)}:} انطلاق العربة بتسارع عالٍ في أجزاء من الثانية حتى تتطابق سرعتها مع سرعة السير الناقل.
    \item \textbf{القص \textenglish{(Cutting)}:} تفعيل آلية القطع أثناء التطابق الحركي، مع إبقاء السكين للأسفل لفترة زمنية محددة لضمان اختراق المنتج.
    \item \textbf{العودة السريعة \textenglish{(Fast Return)}:} رفع الشفرة ورجوع العربة بأقصى سرعة ميكانيكية ممكنة إلى نقطة البداية لانتظار القطعة التالية.
\end{enumerate}

من خلال هذا التوصيف، يتضح أن التحدي الأكبر في هذه المنظومة ليس في عملية القطع بحد ذاتها، بل في \textbf{الخوارزمية البرمجية} القادرة على ربط المتغيرات الديناميكية (السرعة، المسافة، الزمن) ومعالجتها بالتوازي لضمان استقرار الآلة وجودة المنتج النهائي.